\chapter{Normality and Ultraweak Continuity for Positive Functionals}

In what follows, let $M$ be a possibly non-unital $C^{*}$--algebra with its compatible ordered
star-ring structure, and fix a specified Banach predual $P$ of $M$. The existence of $P$ forces
$M$ to be unital, but the public statements below do not assume unitality separately and introduce
the resulting unit only inside their proofs. No nontriviality assumption on $M$ is needed. Let
$\mathcal{P}(M)$ denote the projection lattice of $M$, and let $\varphi$ be a positive
linear functional on $M$. We say that $\varphi$ is \textbf{normal on projections} if the map
$p\mapsto\varphi(p)$ preserves least upper bounds of nonempty directed families of projections.
Equivalently, whenever $(p_{\alpha})$ is an increasing net of projections with supremum $p$, we
have $\varphi(p_{\alpha})\to\varphi(p)$. The formal predicate below is deliberately purely
order-theoretic; it has no topology, $C^{*}$--algebra, or predual built into its definition.

We prove that normality on projections is equivalent to representation by the specified predual,
or equivalently to $\sigma(M,P)$--continuity. Recall that a linear functional $\varphi$ is
\textbf{positive} if $\varphi(x)\ge 0$ whenever $x\ge 0$.

\begin{definition}
  \label{def:normal_on_projections}
  \lean{PositiveLinearMap.IsNormalOnProjections}
  \leanok
  For a positive linear map $f$, \textbf{normality on projections} means Scott continuity of the
  induced map $p\mapsto f(p)$ on the subtype of star projections. Thus $f$ preserves least upper
  bounds of nonempty directed families of projections. The formal definition is stated for
  ordered modules over a semiring and does not assume any topology or $C^{*}$--algebra structure.
\end{definition}

\begin{definition}
  \label{def:specified_ultraweak_dual}
  \lean{Ultraweak.continuousDual,Ultraweak.mem_continuousDual_iff_exists_comp_toUltraweakL}
  \leanok
  \uses{def:sigma_top}
  The \textbf{continuous dual represented by $P$} is the norm-closed subspace of the norm dual of
  $M$ obtained as the image of the canonical isometric embedding of $P$. A norm-continuous
  functional belongs to this subspace exactly when it factors through the explicit continuous
  linear map from $M$ to $\sigma(M,P)$.
\end{definition}

\section{Ultraweakly Continuous Implies Normal}

In what follows, let $T$ denote the set of all $\sigma(M,P)$--continuous positive
linear functionals on $M$, and $E$ the linear space of all finite linear combinations
of elements of $T$. Write $M_{+}$ for the set of positive elements of $M$.

\begin{lemma}
  \label{lem:homeo_of_cts_bij_cpt_haus}
  \lean{isHomeomorph_iff_continuous_bijective}
  \mathlibok
  Every continuous bijection from a compact space to a Hausdorff space is a 
  homeomorphism.
\end{lemma}

\begin{proof}
  \leanok
  Omitted. In Mathlib.
\end{proof}

\begin{theorem}
  \label{thm:Banach_Alaoglu}
  \lean{WeakDual.isCompact_closedBall}
  \mathlibok
  (Banach-Alaoglu) The closed unit ball of the dual space of a normed vector space is compact in the
  weak-$*$ topology.
\end{theorem}

\begin{proof}
  \leanok
  Omitted. In Mathlib.
\end{proof}

\begin{theorem}
  \label{lem:sigma_cont_of_normal_Sak_1_7_4} (Sak 1.7.4)
  \lean{Ultraweak.DirectedOn.exists_isLUB_of_isBounded,Ultraweak.DirectedOn.isLUB_star_right_conjugate}
  \leanok
  \uses{thm:Banach_Alaoglu,lem:homeo_of_cts_bij_cpt_haus,lem:uw_pos_sep_pts,lem:selfadjoint_le_norm,def:sigma_top}
  Every uniformly bounded, increasing net in $M^{s}$ converges
  to its least upper bound in the $\sigma$--topology. Further, if
  $x=\sup_{\lambda}x_{\lambda}$ then $a^{*}xa=\sup_{\lambda}a^{*}x_{\lambda}a$.
\end{theorem}

\begin{proof}
  \leanok
  Letting $T$ denote the set of $\sigma$--continuous positive functionals on $M$ and $E$ its linear span, by
  Lemma \ref{lem:uw_pos_sep_pts}, the subspace $E$ separates the points of $M$ and so the $\sigma(M,E)$
  is Hausdorff. By the Banach-Alaoglu Theorem (\ref{thm:Banach_Alaoglu}) the closed unit ball $B$ of $M$ is
  compact in the $\sigma(M,P)$--topology. Since $\sigma(M,E)$ is weaker than $\sigma(M,P)$, the identity
  map on $B$ is $\sigma(M,P)-\sigma(M,E)$--continuous, and therefore a homeomorphism,
  by Lemma \ref{lem:homeo_of_cts_bij_cpt_haus}. So, to prove that a uniformly bounded net $(x_{t})$ is $\sigma$--
  Cauchy, it's enough to check that for all $\varepsilon>0$ and $\varphi\in T$ that there is a $\Lambda$ such that
  $\varphi(x_{\lambda}-x_{\mu})\le \varepsilon$ for $\lambda \ge \mu \ge \Lambda$.
  
  If $(x_{\lambda})$ is a uniformly bounded increasing net in $M^{s}$, and $\varphi\in T$, then $(\varphi(x_{\lambda}))$ is
  a uniformly bounded increasing net of real numbers and therefore is Cauchy. By the above paragraph, $(x_{\lambda})$ is
  $\sigma$--Cauchy and by the $\sigma$--compactness of $B$ and the fact that $(x_{\lambda})$ is increasing there exists $x\in B$
  so that $x_{\lambda} \to x$ in the $\sigma$--topology. Since for any $\varphi\in T$ and any $\lambda$ we have 
  $0 \le \varphi(x-x_{\lambda})$, by Lemma \ref{lem:non_pos_elem_neg_for_some_state_Sak_1_7_2} we have $0\le x - x_{\lambda}$.
  Therefore $\sup(x_{\lambda})\le x$. For a given $\varphi\in T$ and $\varepsilon = \varphi(x-\sup(x_{\lambda}))$ one find 
  $x_{\lambda}$ so that $\varphi(x-x_{\lambda})<\varphi (x - \sup(x_{\lambda}))$, which is a contradiction, therefore $x=\sup(x_{\lambda})$.

  By Lemma \ref{lem:star_conj_pos} and the fact that $(u^{-1})^{*}=(u^{*})^{-1}$, then $u^{*}x_{\lambda}u\le u^{*}xu$
  is equivalent to $x_{\lambda}\le x$ and so $\sup(u^{*}x_{\lambda}u)=u^{*}\sup(x_{\lambda})u=u^{*}xu$ if $u\in M$ is invertible.

  If $a\in M$ is not invertible, one can shift the spectrum horizontally to avoid zero, i.e. there is a $c>0$ so that $c1+a$ is invertible. Indeed
  observe that $a-t1=(a+c1)-(t+c)1$ and therefore $t\in \sigma(a)\iff t+c \in \sigma(a+c1)$. In particular if we let $t=-c$ here we find that
  $0\in \sigma(a)+c \iff 0 \in \sigma(a+c1)$.
  By the invertible case above, 
  \begin{align}
    c^2\varphi(x_{\lambda})+c\varphi(a^{*}x_{\lambda}+x_{\lambda}a)+\varphi(a^{*}x_{\lambda}a)&=\varphi((c1+a)^{*}x_{\lambda}(c1+a))\\\nonumber
    &\to \varphi((c1+a)^{*}x(c1+a))
  \end{align}
  for all $\varphi\in T$. 
  It suffices to show that $\varphi(a^{*}x_{\lambda}+x_{\lambda}a)\to \varphi(a^{*}x+xa)$ for all $\varphi\in T$, 
  because then $\varphi(a^{*}x_{\lambda}a)\to \varphi(a^{*}xa)$ by the above. We can use Cauchy--Schwarz for this as follows. If $\alpha\le \beta$,
  we have $x_{\beta}-x_{\alpha}\ge 0$, and since the net $(x_{\lambda})$ is uniformly bounded there exists $M>0$ such that (using \ref{lem:selfadjoint_le_norm})
  \begin{align}
    |\varphi(a^{*}(x_{\beta}-x_{\alpha}))|&=|\varphi(a^{*}(x_{\beta}-x_{\alpha})^{1/2}(x_{\beta}-x_{\alpha})^{1/2})|\\\nonumber
                                          &\le \varphi(a^{*}(x_{\beta}-x_{\alpha})a)^{1/2}\varphi((x_{\beta}-x_{\alpha}))^{1/2}\\\nonumber
                                          &\le 2M\varphi(a^{*}a)^{1/2}\varphi((x_{\beta}-x_{\alpha}))^{1/2}
  \end{align}
  and by the positivity of $\varphi$ we have $|\varphi(a^{*}(x_{\beta}-x_{\alpha}))|=|\varphi((x_{\beta}-x_{\alpha})a)|$.
  It follows that $\varphi(a^{*}x_{\lambda}+x_{\lambda}a)\to \varphi(a^{*}x+xa)$, and the proof is finished.
\end{proof}

\begin{theorem}
  \label{thm:uw_continuous_is_normal}
  \lean{PositiveLinearMap.isNormalOnProjections_of_mem_continuousDual}
  \leanok
  \uses{def:normal_on_projections,def:specified_ultraweak_dual,lem:sigma_cont_of_normal_Sak_1_7_4,prop:proj_compl_lat_wstar_Sak_1_10_2}
  If a positive functional $\varphi$ is represented by the specified predual $P$, equivalently if
  $\varphi$ is $\sigma(M,P)$--continuous, then $\varphi$ is normal on projections.
\end{theorem}

\begin{proof}
  \leanok
  Factor $\varphi$ through the explicit map from $M$ to its $\sigma(M,P)$ topology. A nonempty
  directed family of projections converges ultraweakly to any specified least upper bound, so
  continuity of the factor gives convergence of the corresponding values of $\varphi$.
  Positivity makes these values monotone, and hence identifies their limit as their least upper
  bound. Thus the induced map on projections preserves every nonempty directed supremum.
\end{proof}

\section{Normal Implies Ultraweakly Continuous}

If $\mathscr{Q}:\mathcal{P}(M)\to \operatorname{Prop}$ is a predicate, the usual ordering
``$\le$'' on projections induces an order on
$\{p\in \mathcal{P}(M)\mid\mathscr{Q}(p)\}$. We use the same symbol for this induced order.

\begin{lemma}
  \label{lem:exists_uw_ge_normal}
  \lean{Ultraweak.exists_positiveCLM_apply_gt}
  \leanok
  \uses{lem:uw_pos_sep_pts}
  For every positive linear functional $\varphi$ and every nonzero positive element $a\in M_{+}$,
  there exists a positive $\sigma(M,P)$--continuous linear functional $\psi$ such that
  $\varphi(a)<\psi(a)$.
\end{lemma}

\begin{proof}
  \leanok
  By Lemma \ref{lem:uw_pos_sep_pts} there is a $\sigma(M,P)$--continuous positive linear
  functional $\psi_0$ on $M$ such that $\psi_0(a)\ne 0$. Positivity forces
  $\psi_0(a)>0$, so a sufficiently large positive integral multiple of $\psi_0$ strictly
  dominates $\varphi$ at $a$. Neither normality of $\varphi$ nor the assumption that $a$ is a
  projection is needed.
\end{proof}

\begin{lemma}
  \label{lem:msr_th_lemma}
  \lean{PositiveLinearMap.IsNormalOnProjections.exists_nonzero_subprojection_lt_of_chain_lubs,PositiveLinearMap.IsNormalOnProjections.exists_nonzero_subprojection_lt,PositiveLinearMap.IsNormalOnProjections.exists_nonzero_subprojection_lt_of_predual}
  \leanok
  \uses{def:normal_on_projections,prop:proj_compl_lat_wstar_Sak_1_10_2,cor:proj_sub_of_subproj}
  Let $\varphi$ and $\psi$ be positive linear functionals which are both normal on projections.
  If $p\in\mathcal{P}(M)$ satisfies $\varphi(p)<\psi(p)$, then there is a nonzero
  $p_1\in\mathcal{P}(M)$ with $p_1\le p$ such that
  $\varphi(q)<\psi(q)$ for every nonzero projection $q\le p_1$.
\end{lemma}

\begin{proof}
  \leanok
  The reusable formal core assumes only that each nonempty chain of projections has a least upper
  bound. A directed-complete wrapper preserves the stronger order-theoretic interface, and the
  displayed predual wrapper obtains completeness from $P$.

  Let $S$ consist of the subprojections $q\le p$ for which
  $\psi(q)\le\varphi(q)$. It contains zero. The supremum of a chain in $S$ remains below $p$;
  normality of both functionals says that their values at this supremum are the corresponding
  suprema of their values on the chain, so the inequality is preserved. Zorn's lemma therefore
  gives a maximal $q_0\in S$. The strict hypothesis at $p$ implies $q_0\ne p$.

  Put $p_1=p-q_0$. This is a nonzero subprojection of $p$. If a nonzero projection
  $q\le p_1$ failed the desired strict inequality, then $\psi(q)\le\varphi(q)$. The projections
  $q_0$ and $q$ are orthogonal, so $q_0+q$ would be a strictly larger member of $S$, contradicting
  maximality.
\end{proof}

This two-normal-functional formulation is stronger than the specialization needed below. In that
application $\psi$ is initially known to be $\sigma(M,P)$--continuous; Theorem
\ref{thm:uw_continuous_is_normal} first makes it normal, after which this lemma applies.

The following is the key use of normality in the later proof of Theorem
\ref{thm:sigma_cts_of_normal_Sak_1_13_2}: it turns convergence on projections into operator-norm
closure of the corresponding cut-off functionals.

\begin{lemma}
  \label{lem:zorn_base}
  \lean{PositiveLinearMap.IsNormalOnProjections.isUltraweakCutoff_of_isLUB,PositiveLinearMap.IsNormalOnProjections.exists_maximal_isUltraweakCutoff}
  \leanok
  \uses{def:normal_on_projections,def:specified_ultraweak_dual,cor:proj_sub_of_subproj,prop:proj_compl_lat_wstar_Sak_1_10_2}
  Let $\varphi$ be a positive functional which is normal on projections, and let
  $\mathscr{Q}(p)$ mean that the cutoff $x\mapsto\varphi(xp)$ is represented by the specified
  predual $P$. If $S$ is a nonempty directed family of projections with least upper bound $p$ and
  $\mathscr{Q}(q)$ holds for every $q\in S$, then $\mathscr{Q}(p)$ holds. Consequently there is a
  maximal projection $p_0$ satisfying $\mathscr{Q}(p_0)$.
\end{lemma}

\begin{proof}
  \leanok
  For comparable projections $q\le p$, Cauchy--Schwarz gives the operator-norm estimate
  \[
    \|\varphi(\mathord\cdot(p-q))\|
      \le \sqrt{\|\varphi\|}\,\sqrt{\|\varphi(p-q)\|}.
  \]
  Normality identifies $\varphi(p)$ as the least upper bound of the values $\varphi(q)$ for
  $q\in S$. Hence the canonical net over $S$ satisfies $\varphi(p-q)\to0$, and the estimate shows
  that the cutoffs at $q$ converge in operator norm to the cutoff at $p$. The functionals
  represented by $P$ form the norm-closed submodule
  \texttt{Ultraweak.continuousDual}, so the limiting cutoff is represented by $P$ as well.

  The projection lattice supplied by the specified predual is complete, zero has zero cutoff,
  and the preceding directed-supremum closure supplies the chain upper bounds required by Zorn's
  lemma. This yields a maximal projection with represented cutoff.
\end{proof}

\begin{theorem}
  \label{thm:sigma_cts_of_normal_Sak_1_13_2} (Sak 1.13.2)
  \lean{PositiveLinearMap.isNormalOnProjections_iff_mem_continuousDual}
  \leanok
  \uses{def:normal_on_projections,def:specified_ultraweak_dual,thm:uw_continuous_is_normal,lem:zorn_base,lem:exists_uw_ge_normal,lem:msr_th_lemma,lem:cut_down_sigma_closed_cts,thm:real_rank_zero_of_predual,cor:positive_functional_le_of_projection_le,lem:bdd_sigma_cts_iff_bdd_s_cts_Sak_1_8_10}
  For every positive linear functional $\varphi$ on $M$, the following are equivalent:
  \begin{enumerate}
    \item $\varphi$ is normal on projections;
    \item $\varphi$ is represented by the specified predual $P$, equivalently
      $\varphi$ is $\sigma(M,P)$--continuous.
  \end{enumerate}
\end{theorem}

\begin{proof}
  \leanok
  The implication from (2) to (1) is Theorem \ref{thm:uw_continuous_is_normal}. For the converse,
  suppose that $\varphi$ is normal. Lemma \ref{lem:zorn_base} gives a maximal projection
  $p_0\in\mathcal{P}(M)$ for which the cutoff $x\mapsto\varphi(xp_0)$ is represented by $P$.
  The specified predual first supplies the multiplicative identity used in the rest of the proof.
  Suppose that $p_0\ne1$. Lemma \ref{lem:exists_uw_ge_normal} gives a positive
  $\sigma(M,P)$--continuous functional $\psi$ such that
  $\varphi(1-p_0)<\psi(1-p_0)$. By the easy implication, $\psi$ is normal on projections, so the
  two-normal-functional selection Lemma \ref{lem:msr_th_lemma} gives a nonzero projection
  $p\le1-p_0$ such that $\varphi(q)<\psi(q)$ for every nonzero projection $q\le p$.

  Use the explicit analytic corner associated to $p$. Its ambient range is ultraweakly closed, so
  the quotient of $P$ by the preannihilator of that range is a specified Banach predual of the
  corner. Theorem \ref{thm:real_rank_zero_of_predual} therefore makes the corner real rank zero.
  Restrict $\varphi$ and $\psi$ along its explicit inclusion into $M$. Their inequality holds on
  every corner projection: it is strict for nonzero projections and equality at zero. Corollary
  \ref{cor:positive_functional_le_of_projection_le} promotes this to
  $\varphi(a)\le\psi(a)$ for every nonnegative corner element $a$. In particular,
  $px^{*}xp$ is positive in the corner and
  $\varphi(px^{*}xp)\le\psi(px^{*}xp)$. Cauchy--Schwarz now gives
  \begin{align}
    |\varphi(x(p_0+p))|&\le |\varphi(xp_0)|+|\varphi(xp)|\\ \nonumber
    &\le |\varphi(xp_0)|+\varphi(1)^{1/2}\varphi(px^{*}xp)^{1/2}\\ \nonumber
    &\le |\varphi(xp_0)|+\varphi(1)^{1/2}\psi(px^{*}xp)^{1/2}.
  \end{align}
  The cutdown by $p$ is explicitly ultraweakly continuous, so
  $x\mapsto\psi(px^{*}xp)^{1/2}$ is a defining seminorm for the strong topology. The cutoff at
  $p_0$ is ultraweakly, hence strongly, continuous. The estimate therefore makes the cutoff at
  $p_0+p$ strongly continuous; Lemma
  \ref{lem:bdd_sigma_cts_iff_bdd_s_cts_Sak_1_8_10} makes it ultraweakly continuous. Since
  $p\le1-p_0$, the projections $p_0$ and $p$ are orthogonal, so $p_0+p$ is a projection strictly
  above $p_0$. This contradicts maximality. Thus $p_0=1$, and the cutoff at $1$ is precisely
  $\varphi$, proving that $\varphi$ is represented by $P$.
\end{proof}

The theorem is parameterized by an arbitrary specified Banach predual. Consequently, the
ultraweakly continuous positive functionals are determined entirely by the projection order and
are independent of which predual is specified.

\section{Uniqueness of the Predual}

We now let $P$ and $Q$ be two specified Banach preduals of the same possibly non-unital
$C^{*}$--algebra $M$. Either predual supplies unitality internally. Each predual has a canonical
linear isometric embedding into the norm dual $M^{*}$, and its image is the closed subspace of
functionals represented by that predual.

\begin{theorem}
  \label{thm:unique_predual}
  \lean{Ultraweak.continuousDual_eq,Predual.equiv,Predual.equiv_apply_duality,Predual.equiv_eq_of_apply_duality_eq}
  \leanok
  \uses{def:specified_ultraweak_dual,thm:sigma_cts_of_normal_Sak_1_13_2,thm:polar_decomposition_ultraweak_functional,lem:star_left_mul_right_mul_sig_cts_Sak_1_7_8}
  The closed subspaces of $M^{*}$ represented by $P$ and $Q$ are equal. Consequently there is a
  canonical complex-linear isometric equivalence
  \[
    E_{P,Q}:P\mathrel{\cong}Q
  \]
  which preserves evaluation on $M$:
  \[
    \langle x,E_{P,Q}(p)\rangle_{Q}=\langle x,p\rangle_{P}
    \qquad (x\in M,\ p\in P).
  \]
  This pairing-preservation property uniquely characterizes $E_{P,Q}$ among complex-linear
  isometric equivalences $P\cong Q$.
\end{theorem}

\begin{proof}
  \leanok
  Let $f\in M^{*}$ be represented by $P$, and factor it through $\sigma(M,P)$. Split the resulting
  complex-linear functional into its self-adjoint real and imaginary parts. The polar
  decomposition of a self-adjoint ultraweakly continuous functional writes each part as
  $\varphi\circ L_{u}$, where $\varphi$ is positive and $\sigma(M,P)$--continuous and $u$ is a
  self-adjoint unitary. By Theorem \ref{thm:sigma_cts_of_normal_Sak_1_13_2}, normality of the
  positive functional $\varphi$ is intrinsic to the projection order, so $\varphi$ is also
  $\sigma(M,Q)$--continuous. Left multiplication by $u$ is $\sigma(M,Q)$--continuous by Lemma
  \ref{lem:star_left_mul_right_mul_sig_cts_Sak_1_7_8}. Thus both self-adjoint parts, and hence
  $f$, are represented by $Q$. Interchanging $P$ and $Q$ proves equality of the two represented
  closed subspaces of $M^{*}$.

  The canonical embeddings identify both $P$ and $Q$ linearly and isometrically with this common
  subspace. Compose the first identification with the inverse of the second, using the identity on
  the common image, to obtain $E_{P,Q}$. By construction its image represents exactly the same
  functional on $M$, which proves the evaluation formula. Finally, the canonical embedding of
  $Q$ into $M^{*}$ is injective, so that formula determines the value of any pairing-preserving
  equivalence at every $p\in P$ and proves uniqueness.
\end{proof}
